\section{Horizontal links extracted from a long geodesic}
\label{links:section}

The aim is to bound the capped mean cost of a local horizontal crossing.
A long geodesic toward an endpoint of the putative flat has negligible
raw excess.  We extract many crossings from it, control the cost lost
at exceptional pieces by the charged-tunnel estimates, and use row
concentration to compare the surviving sum with a single-link capped mean.
The last step requires two truncation estimates, proved separately below.
Only extremality of the endpoint is used in this extraction.

Costs remain in lattice units, while positions belong to
$\mathcal G_n=n^{-1}\mathbb Z^2$.  Put $d=1-s^2$, and use coordinates $(u,y)$
in the basis $(v,B)$; thus $A=(d,s)$ and $B=(0,1)$.
Tracking and radius-$10$ locality are still measured in the fixed
$(v,w)$ maximum norm.  For $s<1/4$, changes to either
$(v,B)$ or $(A,B)$ coordinates and back cost a factor at most
$2$.  Thus a $B$-wall stays within $2r$ of its nominal line
in the present coordinates.  Lattice-edge rounding errors are at
most $2K_{\rm lat}/n$ in this norm and are absorbed by taking $n$ large.

Constants may depend on the fixed flat.  Those subscripted by
$s$ or $K$, and the large constants used only in the proof, may
also depend on the indicated auxiliary parameters.  The dependence
radius of a link is fixed independently of all truncation and
coarse-graining choices.

\subsection{The signed link and its translation conventions}

Let $W_j$ be the $B$-road with nominal line
$\{(jd,js+t):t\in\mathbb R\}$.  We keep its ordered occurrence
labels, so two visits to one lattice vertex remain distinct.
The nominal parameter is continuous and nondecreasing, and a fixed
first-occurrence convention selects its marked cuts.  This convention
also selects a unique cut at a pause, where the occurrences retain
their order.  The road is locally finite.

Write $R_j(t,t')$ for the signed raw $Q_n$-cost between the
indicated occurrences at $t$ and $t'$.  Reversing the arguments
negates this bookkeeping quantity; a physical reverse traversal has
its own positive passage cost.  With the interpolated local gauge
$g_j$, put
\[
 C_j(t,t')=R_j(t,t')-g_j(t')+g_j(t).
\]
Signed costs are additive in occurrence labels.  Interpolation gives
\[
 C_j(t,t')=(t'-t)C_j(\kappa,\kappa+1),
 \qquad \kappa=\phi_j+m,\quad m\in\mathbb Z,
\]
for occurrences in the marked wall period
$[\kappa,\kappa+1]$, with the relative wall phase $\phi_j$ from
Section~\ref{phases:section}.  This includes pauses.  Longer
intervals are split at the actual marked levels.

Fix a rational row spacing $\sigma<s/1000$, chosen so that $s$ and
all later row offsets are integral multiples of $\sigma$.  Write
$y_k=k\sigma$, $b=s/50$ and $b'=s/20$.  Let $c_{j,k}$ be the
center occurrence of $W_j$ with nominal height $y_k$, that is with
wall parameter $t=y_k-js$, selected by the first-occurrence convention.
An admissible crossing from $W_j$ to $W_{j+1}$ at row $k$ is a finite
physical path $\gamma$, with endpoint occurrences $x\in W_j$ and
$x'\in W_{j+1}$, such that
\begin{enumerate}
\item every point has $|y-y_k|\le b$ and
$jd-b'\le u\le(j+1)d+b'$;
\item for any two points in their traversal order, their horizontal
coordinates satisfy $u_2\ge u_1-b'$.
\end{enumerate}
The second condition bounds every horizontal drawdown along the
crossing, including one that occurs after an earlier visit to the
terminal wall.  At finite $n$, round the band boundaries outward to
the nearest lattice boundaries; the $O(1/n)$ error fits inside the
strict margins used below.  Equivalently, fix rational padded
boundaries and require sufficient divisibility of $n$.
Measure the crossing at the two row centers by setting
\begin{align}
 L_{j,k}(\gamma)
  &=C_j(c_{j,k},x)
    +Q_n(\gamma)-g_{j+1}(x')+g_j(x)
    +C_{j+1}(x',c_{j+1,k}), \\
 X_{j,k}&=\min_\gamma L_{j,k}(\gamma).
 \label{links:definition}
\end{align}
The closed-window endpoint-selection argument preceding
Lemma~\ref{found:lem-envelope} gives an attained measurable local
minimum. The band and all-pairs drawdown conditions are closed and
survive loop deletion. The class is nonempty for large $n$: follow the horizontal center line by a
lattice path and retain the portion between appropriate consecutive
crossings of the two proper vertical roads.  The planar
rectangle-crossing property ensures these intersections even when
the roads have self-intersections.

When a link is used in a new road, its minimizer may be chosen with
vertex endpoints. For a loop-erased minimizer starting at an interior
point $x$ of a wall edge, let $v$ be the first adjacent vertex reached
by the crossing. Delete the partial prefix from $x$ to $v$ and move
the wall attachment to the corresponding occurrence of $v$. The change
in the center-adjusted cost is
\[
 R_j(x,v)-Q_n(x\mathbin{\to}v)
 =\begin{cases}
 0,&\text{if the partial traversal follows the wall},\\
 -2\tau(x\mathbin{\to}v),&\text{if it reverses the wall}.
 \end{cases}
\]
The right endpoint has the same property after deleting its final
partial edge. The two walls have positive separation, so for large
$n$ these deletions are disjoint. The remaining subpath retains the
band and drawdown conditions. Its cost cannot exceed the minimum and
therefore equals it. Choose the first such vertex-ended minimizer in
the relative order above. This supplies a lattice path at every link
splice, while the definition still compares all allowed interior
attachments.

The minimum has a uniform bound on its negative first moment:
Lemma~\ref{found:lem-envelope} gives
\begin{equation}
 \mathbb E(X_{j,k})^-\le C\beta.
 \label{links:negative-moment}
\end{equation}
To see why taking a minimum causes no loss, apply calibration
domination to each physical $\gamma$.  After taking the difference
between the two center labels, the endpoint calibration--gauge
discrepancies have $L^1$ suprema at most $C\beta$.  The two
interpolated center arms obey the same bound.  Their suprema provide
one lower bound valid for all competitors before minimization.
The calibration enters this estimate; the definition of $X_{j,k}$
uses only local inputs.  The constant $C$ is independent of small
$s$, the row mesh, and the input rule, because only two labels and
bounded parameter windows are compared, and finite averaging of
gauges preserves these bounds.

The same first-moment proof applies with every phase entering the link
and its gauge projections fixed. The calibration depends only on the
edges, while the local-envelope proof treats phases as parameters and
uses period, interpolation, and projection bounds uniform over their
admissible assignments. Thus a link \(X^t\) evaluated under any fixed
local phase template of Section~\ref{access:phases} satisfies
\begin{equation}
 \mathbb E[(X^t)^-]\le C\beta.
 \label{links:fixed-phase-negative-moment}
\end{equation}
This fixed-phase form will be used for the sparse negative-link flags.

Two properties of the phase construction will turn selected sums
into a statement about one link.  Use Lemma~\ref{phases:locality}
with $\mathcal H$ containing $s$, $\sigma$, and all required row
offsets.  Translation by $d v=A-sB$ advances a column at fixed
height; translation by $\sigma B$ advances a row.  Covariance of
the phases and averaged gauges therefore gives every $X_{j,k}$
the same law, denoted by $X$.  The lemma also gives a fixed integer
$q$ for which entire-column sigma-fields with indices in one residue
class modulo $q$ are mutually independent.  These are precisely
the stationarity and independence used in Lemma~\ref{probability:rows}.
For a finite family of crossing templates, apply this construction
and the extraction separately to each template.

\subsection{The geometric extraction}

Two geometric facts are needed. Near-minimal norm length forces most
fine increments toward the extreme endpoint. At the coarser scale,
deleting neighborhoods of bad pieces gives the all-pairs order required
by the tunnel-tail estimate. Write
\[
 \mathcal C=\{z:\mu(z)=h(z)\}
\]
for the cone on the flat, and choose pointed closed cones with
\[
 \mathcal C\setminus\{0\}\subset\operatorname{int}\mathcal C_1,
 \qquad
 \mathcal C_1\setminus\{0\}\subset\operatorname{int}\mathcal C_2.
\]
They can be arbitrarily close to \(\mathcal C\). On every such cone
sufficiently close to \(\mathcal C\), \(h(z)\ge c|z|\), with \(c>0\)
independent of the enlargement.

\begin{lemma}[A uniform modulus at an extreme point]
\label{links:extreme}
Let \(a\) be an extreme point of the norm ball, let
\(S_\mu=\{z:\mu(z)=1\}\), and let \(U\) be a relatively open
neighborhood of \(a\) in \(S_\mu\). For \(t\ge0\), put
\[
 \omega_U(t)=\sup\left\{\nu(S_\mu\setminus U):
    \nu\in\mathcal P(S_\mu),\
    \left|\int z\,d\nu(z)-a\right|\le t\right\}.
\]
Then \(\omega_U(t)\to0\) as \(t\downarrow0\). If \(N>0\),
\(0\le d_0<1/2\), and finitely many increments satisfy
\[
 \left|\sum_i z_i-Na\right|\le d_0N,\qquad
 \left|\sum_i\mu(z_i)-N\right|\le d_0N,
\]
then, with \(C_a=2(1+|a|)\),
\begin{equation}
 \sum_{i:z_i/\mu(z_i)\notin U}\mu(z_i)
 \le (1+d_0)N\,\omega_U(C_a d_0).
 \label{links:extreme-modulus}
\end{equation}
Zero increments are omitted. In particular, endpoint error \(o(N)\)
and total norm length \(N+o(N)\) imply \(o(N)\) norm mass outside \(U\).
\end{lemma}

\begin{proof}
Put \(F=S_\mu\setminus U\). If the first assertion failed, there
would be \(t_j\downarrow0\) and probability measures \(\nu_j\) with
barycenters within \(t_j\) of \(a\) and \(\nu_j(F)\ge e>0\).
A weakly convergent subsequence exists on the compact sphere. Its
limit \(\nu\) has barycenter \(a\), and the closed-set inequality
gives \(\nu(F)\ge\limsup_j\nu_j(F)\ge e\).

The convex hull of \(F\) is compact and omits \(a\). Indeed,
Carath\'eodory's theorem would otherwise express \(a\) as a finite
nontrivial convex combination of points different from \(a\),
contrary to extremality. If \(\nu(F)=1\), its barycenter lies in that
convex hull. If \(0<\nu(F)<1\), the conditional barycenters on \(F\)
and its complement lie in the norm ball and express \(a\) as a
nontrivial convex combination. Extremality makes both equal to \(a\),
whereas the first lies in the convex hull of \(F\). Both cases are
impossible, proving \(\omega_U(t)\to0\).

For the second assertion put \(S=\sum_i\mu(z_i)>0\) and give
\(z_i/\mu(z_i)\) mass \(\mu(z_i)/S\). Its barycenter is
\((\sum_i z_i)/S\), at distance at most
\[
 \frac{d_0+d_0|a|}{1-d_0}\le C_a d_0
\]
from \(a\). The mass on \(F\), multiplied by \(S\le(1+d_0)N\),
is the left side of \eqref{links:extreme-modulus}.
\end{proof}

Deleting only the bad indices would not control an interval that
crosses several of them. The next lemma removes a density halo:
every interval touching a retained index then contains enough good
increments to enforce forward order. The raw penalty on bad indices
pays for all removed effective gauge differences.

\begin{lemma}[Removal of bad intervals]
\label{links:halos}
Let \(x_0,\ldots,x_m\) be successive cuts of a path. All distances
in this lemma use the \((A,B)\) maximum norm. Suppose
\(K/2\le|x_{i+1}-x_i|\le CK\), and associate to each starting cut
\(x_i\) a center \(z_i\) within a fixed distance \(H_0\), independent
of the cone enlargement. Declare an index bad when its
increment is outside \(\mathcal C_1\), or its physical subpath has
a backward \(h\)-excursion of size at least one. Suppose, for
\(c_*>0\) and \(\epsilon_n^{\rm raw}\to0\), its raw costs satisfy
\[
 Q_i\ge-\epsilon_n^{\rm raw}n\quad\hbox{for all }i,
 \qquad Q_i\ge c_*n\quad\hbox{for bad }i.
\]
Let effective cut labels satisfy
\(|\widetilde g_i-g_{\rm ref}|\le\lambda n\) for one common constant
\(g_{\rm ref}\). For sufficiently small \(\lambda>0\) there is a set
\(R\) of removed indices containing all bad indices, of cardinality
at most a fixed multiple of their number, for which
\begin{enumerate}
\item \(\sum_{i\in R}(Q_i-\widetilde g_{i+1}+\widetilde g_i)\ge0\);
\item retained centers \(z_i\) are in \(\mathcal C_2\)-order up to a bounded
additive error, for every pair of retained indices;
\item the same order includes the initial and terminal points, up to
a bounded additive error;
\item repetitions of any one center have bounded multiplicity, independently of how
small the cone enlargement is.
\end{enumerate}
\end{lemma}

\begin{proof}
Let \(B_0\) be the bad set. For \(\theta>0\), put
\[
 R=\{i:\text{some index interval }I\ni i
                    \text{ has }|I\cap B_0|>\theta|I|\}.
\]
A one-dimensional maximal-interval bound gives
\(|R|\le3|B_0|/\theta\): select disjoint witness intervals in
decreasing order of length. Their triple enlargements cover all
discarded witness intervals, and their lengths sum to at most
\(|B_0|/\theta\). Every interval containing a retained index
therefore has bad density at most \(\theta\).

Let \(\ell_1,\ell_2\) be the coordinate functionals of
\(\mathcal C_2\). Compactness and the strict cone inclusion give
\(\ell_j(x_{i+1}-x_i)\ge\kappa_0K\) at good indices, whereas at bad
indices they are at least \(-CK\). Choose
\(\theta<\kappa_0/(2(C+\kappa_0))\). Every interval containing a
retained index then satisfies
\[
 \ell_j\left(\sum_{i\in I}(x_{i+1}-x_i)\right)
                      \ge\tfrac12\kappa_0K|I|.
\]
Apply this to the interval between two retained cut indices, then
absorb the bounded errors in choosing their centers. To include the
initial endpoint and a retained starting-cut representative \(x_i\),
use the prefix through increment \(i\) and subtract that increment,
whose size is at most \(CK\). The suffix from \(x_i\) already contains
the retained index. Center rounding adds another bounded error.
The coordinate matrices range in a fixed compact set, so these errors
remain uniform as the cone enlargement shrinks.

For multiplicity, apply the same density argument to \(h\). After
decreasing \(\theta\) if necessary, it gives
\(h(x_j-x_i)\ge cK(j-i)\) on the relevant intervals, where \(c\) is
independent of the cone enlargement. Only boundedly many retained
centers can therefore lie in any fixed-radius ball.

For the cost assertion, bound each removed gauge difference separately:
\begin{align*}
 \sum_{i\in R}(Q_i-\widetilde g_{i+1}+\widetilde g_i)
 &\ge n\{c_*|B_0|-(\epsilon_n^{\rm raw}+2\lambda)|R|\}\\
 &\ge n\left\{c_*-\frac{3(\epsilon_n^{\rm raw}+2\lambda)}{\theta}\right\}|B_0|.
\end{align*}
Choose \(\lambda<c_*\theta/24\) and then \(n\) large. The bracket is
at least \(c_*/2\), retaining a positive penalty per bad index.
The set \(R\) depends only on the bad indices, so this conclusion
holds for every assignment of effective gauges satisfying the stated
bound.
\end{proof}

\paragraph{Costs at shared cuts.}
For an allowed road label \(W\), a cut occurrence \(x\), and a
permitted datum \(\bar x\) on the same wall, define
\begin{equation}
 \widetilde g_W(x;\bar x)=g_W(\bar x)+R_W(\bar x,x).
 \label{links:effective-gauge}
\end{equation}
The datum is a prescribed row center or reference center in the
bounded windows of the access construction. At a shared cut \(x_i\),
both incident pieces will use one datum \(\bar x_i\). Writing
\(\widetilde g_i=\widetilde g_{W_i}(x_i;\bar x_i)\), define
\begin{align}
 c_i^{\rm cost}
 &=\beta^{-1}\{C(\bar x_i,x_i)+Q_n(\gamma_i)
       -g(x_{i+1})+g(x_i)+C(x_{i+1},\bar x_{i+1})\}\notag\\
 &=\beta^{-1}\{Q_n(\gamma_i)-\widetilde g_{i+1}+\widetilde g_i\}.
 \label{links:piece-cost}
\end{align}
The second equality is additivity of signed wall costs. The cut
selection takes its suprema over every allowed label, cut occurrence,
and future datum before any row is chosen. Lemma~\ref{found:lem-gauge-tail}
therefore bounds all the effective gauges in
\eqref{links:effective-gauge} relative to one deterministic reference
mean for the law and scale.

\subsection{Uniform extraction}

We state the output shared by the two capped estimates. For the final
fixed rational directions, let \(\mathfrak R_n\) be the class of
admissible input-road laws with
\[
 c_0\le\beta\le C_s n^\alpha,\qquad \alpha<1,
\]
where these constants are fixed as the law and \(n\) vary. The
preceding local, selected-cut, and collection estimates apply
uniformly to this whole class.

The order of choices is as follows. Fix \(s,\sigma\), the local access
construction, and \(\chi=0.99\). Choose an integer \(L>99\), put
\(N=n^L\), and take \(m=m_*>4\) in Lemma~\ref{access:selected-cuts}
with \(m_*\vartheta>2L+10\). Retain its radii \(H_m\) and \(K>20H_m\);
\(K\) is the actual preliminary exit radius. Fix the finite tunnel
types and their collection constants in one larger cone. The capped
proofs then choose \(F,T\), their caps or finite cap menu, and all
entropy tolerances. For their requested thin width \(\delta\) and
row tolerance \(\eta\), choose the final cones, the halo parameters
\(\theta,\lambda\), the fine cone and radius, and the rational
direction errors. Finally impose divisibility and take \(n\) large.
The constants multiplying \(\eta\) below are fixed before \(\eta\);
later fine and cone choices enter the error sequence. The finite phase
lists are independent of the eventual direction denominators, as
explained in Section~\ref{phases:section}. These proof parameters
never enlarge the input support of an individual link.

\begin{proposition}[Uniform extraction from a long geodesic]
\label{links:extraction}
For every sufficiently small fixed \(\delta,\eta>0\), the final
geometric choices above can be made so that there are deterministic
\(\epsilon_n,\pi_n\to0\) and fixed constants \(H,C,C_K,c>0\) and
integers \(q_K,m_0\) with the following property. For each
\(\mathfrak r\in\mathfrak R_n\) there is an event \(\Omega_{n,\mathfrak r}\)
such that
\begin{equation}
 \sup_{\mathfrak r\in\mathfrak R_n}
 \mathbb P_{\mathfrak r}(\Omega_{n,\mathfrak r}^{\,c})\le\pi_n.
 \label{links:extraction-event}
\end{equation}
Let \(x_n^\star\) be a deterministic nearest vertex of \(\mathcal G_n\)
to \(Na\), and put
\[
 j_n=\left\lfloor u(x_n^\star)/d+\tfrac12\right\rfloor,\qquad
 \mathcal J_n=\{0,1,\ldots,j_n-1\}.
\]
For large \(n\), \(cN\le j_n\le CN\). On \(\Omega_{n,\mathfrak r}\),
a geodesic from \(0\) to \(x_n^\star\) has consecutive selected pieces
\(\gamma_i:x_i\to x_{i+1}\), \(0\le i<m_N\), after two bounded terminal
fragments are removed, and the following objects and bounds.
\begin{enumerate}
\item There are at most \(CN/K+C\) pieces, each of diameter \(O(K)\),
with consecutive cut displacement between \(K/2\) and \(2K\) in the
\((A,B)\) maximum norm. The shared data in
\eqref{links:piece-cost} have
\(|\widetilde g_i-g_{\rm ref}|\le\lambda n\). A removed set \(R\)
contains every bad piece of Lemma~\ref{links:halos} and satisfies
\begin{equation}
 |R|\le\epsilon_nN,\qquad
 \sum_{i\in R}c_i^{\rm cost}\ge0,\qquad
 \sum_{i\notin R}c_i^{\rm cost}\le\epsilon_nN.
 \label{links:retained-cost}
\end{equation}

\item The retained pieces have types \(t_i\) from the fixed access
family and centers \(z_i\) on its fixed finite family of full-rank
center lattices, after bounded rounding. Let \(E=[e_1\ e_2]\), where
the \(e_j\) are the boundary generators of \(\mathcal C_2\), normalized
by \(h(e_j)=1\) and ordered so that \(e_1\to a\) and \(e_2\to b\)
as \(\mathcal C_2\) shrinks to \(\mathcal C\). In
chronological order the centers satisfy
\begin{equation}
 E^{-1}(z_j-z_i)\ge-H(1,1)\quad(i<j),\qquad
 z_i\in E([-H,CN+H]\times[-H,\delta N+H]).
 \label{links:retained-geometry}
\end{equation}
The inequality is coordinatewise. Each rounded site occurs at most
\(m_0\) times, and the initial distance from the origin plus total
successive center distance is at most \(C(N+1)\). These bounds persist
on every type and sparse-residue subsequence. The matrices \(E\) range
in a fixed compact family, and \(H,C,m_0\) remain bounded when the
final cone is narrowed.

\item An irregular retained set \(\mathcal I\) has
\(|\mathcal I|\le(C_K\eta+\epsilon_n)N\). Each
\(i\notin R\cup\mathcal I\) supplies at most \(q_K\) admissible links
\(\mathcal L_i\), with
\begin{equation}
 c_i^{\rm cost}\ge\sum_{(j,k)\in\mathcal L_i}X_{j,k}/\beta.
 \label{links:unclipped-comparison}
\end{equation}
In the chronological list
\(\mathcal L=\bigcup_{i\notin R\cup\mathcal I}\mathcal L_i\), columns
are distinct, increasing, and belong to \(\mathcal J_n\). For
sufficiently small fixed \(\eta\), \(|\mathcal L|\ge cN\) for all
large \(n\), after decreasing \(c\) once. If \(k_1,\ldots,k_\ell\)
are the selected rows, then
\begin{equation}
 \begin{gathered}
 |\mathcal J_n\setminus\{j:(j,k)\in\mathcal L\}|
       \le(C_K\eta+\epsilon_n)N,\qquad |k_1|\le j_n^2,\\
 \sum_{\nu<\ell}|k_{\nu+1}-k_\nu|
       \le(C_K\eta+\epsilon_n)N.
 \end{gathered}
 \label{links:selected-rows}
\end{equation}
Rows can be held across omitted columns without increasing their variation.

\item Every retained piece obeys the charged-tunnel bound
\begin{equation}
 (c_i^{\rm cost})^-\le Y_{z_i,t_i}+E_{z_i}.
 \label{links:tunnel-bound}
\end{equation}
Here \(Y\) is the local saving of Section~\ref{access:tunnels}, and
\(E\) is a fixed local sum at the piece center that dominates both
selected endpoint errors for every allowed piece. The finite
phase-template fields for \(Y\) satisfy
Proposition~\ref{access:collection}; those for \(E\) have fixed range
and tail \(C_K\min(1,t^{-m_*})\).
\end{enumerate}
Set every selected list and count, including \(m_N,R,\mathcal I,\mathcal L\),
to the empty or zero value off \(\Omega_{n,\mathfrak r}\).
For finitely many link templates, apply the proposition separately,
intersect their events, and take the maximum of their error sequences
and the sum of their failure sequences.
\end{proposition}

\begin{proof}
\emph{An edge-only event and fine directions.}
Choose the nearest endpoint and the geodesic by fixed tie conventions.
The metric and cheap-path estimates give an event \(\mathsf M_n\)
depending only on the edge weights, with
\(\rho_n=\mathbb P_{\rm edge}(\mathsf M_n^c)\to0\), on which
\begin{equation}
 |Q_n(\gamma)|\le q_nN,\qquad |\gamma|_{\rm edges}\le CNn,\qquad
 \gamma\subset[-CN,CN]^2,\qquad
 q_n=C\frac{(Nn)^\chi+1}{N}\to0.
 \label{links:long-geodesic}
\end{equation}
Here \(L>99\) makes \((Nn)^{0.99}=o(N)\); endpoint rounding contributes
only \(O(1)\) to \(Q_n\). Write \(T_n(x,y)=T(nx,ny)\) for scaled
metric-graph endpoints. On the same event, every such pair in a fixed
enlargement of this box and at the bounded scaled distances used below
satisfies
\[
 T_n(x,y)\ge n\mu(y-x)-e_n,\qquad e_n=C_Kn^\chi=o(n).
\]
The graph-endpoint assertion includes the maximum incident-edge bound
from Lemma~\ref{model:metric}. Thus it holds simultaneously before any
cuts are chosen.

The fine choices in the stated parameter order can be made as follows.
Choose the rational approximation to \(a\), a neighborhood \(U\), and
\(\varepsilon_f>0\) so that every increment directed in \(U\) satisfies
\[
 u(z)\ge0,\qquad |y(z)|\le\varepsilon_f u(z)
                         \le\tfrac12\eta\sigma\,\mu(z).
\]
Indeed \(u(a)>0\), \(y(a)\to0\) as \(v\to a\), and \(u(z)\le C\mu(z)\)
by fixed norm comparison. Take \(\varepsilon_f\) small enough and choose
the fixed radius \(\varrho\) so that
\begin{equation}
 \kappa=s+d\varepsilon_f\le\tfrac12,\quad
 12\varrho<H_m,\quad \varrho<\sigma/100,\quad
 2(1+s)\varrho<b/10,\quad 2r+4d\varrho<b'.
 \label{links:fine-choices}
\end{equation}
These conditions are compatible because \(s<1/4\) and \(2r<b'\).

Use the fine exit partition of radius \(\varrho\) in the \((A,B)\)
maximum norm. Each fine trace has diameter at most \(2\varrho\),
including the final incomplete trace. The edge budget bounds the number
of fine segments by \(CN/\varrho+1\). For their increments \(z_i\), put
\(S_n=\sum_i\mu(z_i)\). Summing the metric lower bounds gives
\[
 nS_n\le nN+C+q_nN+(CN/\varrho+1)e_n,
\]
while the triangle inequality and endpoint rounding give
\(S_n\ge N-C/n\) and \(|\sum_i z_i-Na|\le C/n\). Thus both errors in
Lemma~\ref{links:extreme}, divided by \(N\), are at most
\[
 d_n=C\left(\frac1{nN}+\frac{q_n}{n}
              +\frac{e_n}{n\varrho}+\frac{e_n}{nN}\right)\to0.
\]
Define
\[
 \epsilon_n^{\rm fine}=(1+d_n)\omega_U(C_a d_n)\to0.
\]
The total norm mass outside \(U\) is at most
\(\epsilon_n^{\rm fine}N\). Every full fine increment has norm length
at least \(c\varrho\), so there are at most
\(C\epsilon_n^{\rm fine}N/\varrho+1\) bad increments, including a
possible bad incomplete one. A zero final segment is omitted.
The good transverse variation is at most
\(\tfrac12\eta\sigma S_n\le\eta\sigma N\) for large \(n\), and the bad
variation is at most \(C\epsilon_n^{\rm fine}N\). All these bounds hold
on the same edge event for every road law.

\emph{Cuts and the uniform good event.}
Place the actual preliminary \(K\)-exit marks in the \((A,B)\) maximum
norm. At each mark test every fine interval whose occurrence-time
interval meets the terminal arc, including at either endpoint.
The coordinate identities
\[
 u=d\xi_1,\qquad y=s\xi_1+\xi_2
\]
put every good fine displacement in \(\mathcal G_\kappa\) of
\eqref{access:horizontal-hyp}. With \(\delta_f=2\varrho\),
\(\kappa\le1/2\) and \(12\varrho<H_m\) imply its exact strict
condition
\[
 H_m>\frac{2(1+\kappa)}{1-\kappa}\delta_f.
\]
If every fine interval meeting the terminal arc is good, use the
\(B\)-wall candidates of Lemma~\ref{access:selected-cuts}; otherwise
use its annulus candidates. Both tuple families and all future datum
choices were included before this test in the single objective
\(W_c/\beta+V_c/a_n\).

With \(a_n=C_1n^{1-\vartheta}\), that lemma supplies
\begin{equation}
 \mathbb P_{\mathfrak r}(E_z>t)\le C_K\min(1,t^{-m_*}),\qquad
 W_{\rm cut}/\beta\le E_z,\qquad V_{\rm cut}\le a_nE_z.
 \label{links:joint-cut-bound}
\end{equation}
The inequalities cover every allowed effective gauge at the chosen
cut. Let \(\mathcal Z_n\) be the deterministic set of possible mark
cells in the box of \eqref{links:long-geodesic}, with its fixed
boundary enlargement. It has at most \(Cn^{2L+10}\) cells. Define
\[
 \Omega_{n,\mathfrak r}
   =\mathsf M_n\cap\{\max_{z\in\mathcal Z_n}a_nE_z\le\lambda n\}.
\]
The union bound in the selected-cuts lemma gives
\[
 \mathbb P_{\mathfrak r}(\Omega_{n,\mathfrak r}^{\,c})
 \le \rho_n+C_{\lambda,K}n^{\,2L+10-m_*\vartheta}
 =:\pi_n\longrightarrow0
\]
uniformly in \(\mathfrak r\). No independence between the metric event
and the cut fields is needed. A selected cut lies within \(H_m\) of its
mark. Since successive preliminary displacements equal \(K\), selected
displacements lie between \(K-2H_m\) and \(K+2H_m\), in particular
between \(K/2\) and \(2K\). The intervening pieces and the two terminal
fragments have diameter \(O(K)\).

\emph{The halo and its cost.}
Each terminal fragment has raw cost at least \(-e_n\), so the selected
pieces have total raw cost at most \(q_nN+2e_n\). Every such piece has
raw cost at least \(-e_n\). For an off-\(\mathcal C_1\) displacement
in the fixed annulus, compactness gives a positive lower bound for
\(\mu(z)-h(z)\). For a unit backward \(h\)-excursion, split before,
during, and after the excursion; \(\mu\ge|h|\) gives raw cost at least
\(2n-3e_n\). For some fixed \(c_*>0\), bad pieces therefore cost at
least \(c_*n\) for large \(n\). Let \(B_0\) be their index set.
The edge budget gives at most \(CN/K+C\) pieces, whence
\[
 \frac{|B_0|}{N}
 \le C\left(\frac{q_n}{n}+\frac{e_n}{nK}
                         +\frac{e_n}{nN}\right).
\]
Use the fixed \(\lambda<c_*\theta/24\) from the halo lemma and
\(\epsilon_n^{\rm raw}=e_n/n\). The halo set \(R\) depends only on
\(B_0\), and \(|R|\le3|B_0|/\theta\). On \(\Omega_{n,\mathfrak r}\)
the gauge bound holds simultaneously for every later datum, so the
removed cost is nonnegative for the eventual shared choices. They
telescope on the full consecutive list, giving
\[
 \sum_{i\notin R}c_i^{\rm cost}
 \le\sum_i c_i^{\rm cost}
 \le c_0^{-1}(q_nN+2e_n+2\lambda n).
\]

The halo's endpoint order puts the retained centers in
\eqref{links:retained-geometry}: the second \(E\)-coordinate of \(a\)
tends to zero as \(\mathcal C_2\) shrinks to \(\mathcal C\).
The endpoint and center errors stay bounded because \(E\) ranges in a
compact family. Each preliminary exit uses at least \(cKn\) edges,
so there are \(O(N/K)\) centers. Their successive-distance budget,
including the initial \(O(K)\) distance, is \(O(N)\). The halo lemma
also gives the stated multiplicity bound.

\emph{Rows and regular links.}
Track the fine polygon by hysteresis: keep its current row until it
leaves the central sub-band of half-width \(b/4\), then choose the
nearest row. Since \(\sigma<b/10\), each change except possibly the
first consumes at least \(b/10\) of transverse variation. The row-index
variation is bounded by a constant times that variation divided by
\(\sigma\), plus a constant times the number of changes. It is at most
\[
 C_K\eta N+C_{\rm fixed}\epsilon_n^{\rm fine}N+O(1).
\]
Let \(\mathcal I_0\) contain exactly the pieces whose closed
occurrence-time intervals meet a row change or a bad fine interval,
or which are incident to a failed terminal test. Define this set on
the full selected list, regardless of membership in \(R\).
A row-change time lies in at most two closed coarse intervals.
A fine trace has diameter
at most \(2\varrho\). Its closed occurrence-time interval can meet at most
two closed coarse intervals, since otherwise it contains two consecutive cuts
separated by \(K/2\). Each terminal arc occurs after its preceding
preliminary mark. A fine interval meeting three terminal arcs would
contain two consecutive preliminary marks, separated by \(K\).
Each failed test is incident to at most two pieces and is charged to a
bad fine interval meeting its terminal arc. Thus
\[
 |\mathcal I_0|
 \le C_K\eta N+C_{\rm fixed}\epsilon_n^{\rm fine}N+O(1).
\]

On a piece outside \(\mathcal I_0\), every fine interval meeting it
chronologically is good. The polygon calculation in
\eqref{access:horizontal-submotion}, now applied to any
submotion of this piece with \(\delta_f=2\varrho\), gives horizontal
drawdown at most \(4d\varrho\). The path is within transverse distance
\(2(1+s)\varrho\) of its fine polygon. The hysteresis and
\eqref{links:fine-choices} therefore keep the whole piece in one full
row band, with drawdown less than \(b'\). Its two successful terminal
tests supply \(B\)-wall vertex occurrences.
Write \(k_i\) for this constant hysteresis row. Two adjacent pieces
outside \(\mathcal I_0\) have the same row: their closed time intervals
share a cut, and a change there would put the incident pieces in
\(\mathcal I_0\).

The first \((A,B)\) coordinate increases across this piece. Its
maximum norm displacement is at least \(K-2H_m\). If the negative
first coordinate realized that norm, its size would be at most
\(4\varrho\) by \eqref{access:horizontal-submotion}; if the second
coordinate realized it, its size would be at most
\(4(1+\kappa)\varrho/(1-\kappa)\le12\varrho\).
Both contradict \(12\varrho<H_m<K/20\). Hence the first-coordinate
increase is at least \(K-2H_m\). Since \(\rho=2r\) is the wall tracking
width in these coordinates, the terminal wall index exceeds the initial
one by at least \(K-2H_m-2\rho>0\).

Split the piece at successive first meetings of the next \(B\)-wall.
At each intermediate meeting, use one occurrence of that wall for both
adjoining subpaths.
Planarity forces a meeting before passage beyond that wall's far
padded side, and the drawdown bound for every submotion supplies the
backward allowance. End the final subpath at the original terminal cut,
even if it met that wall earlier. Its overshoot is at most \(2r+4d\varrho<b'\),
so the entire trailing portion is retained. These are admissible links
in \eqref{links:definition}; the piece diameter bounds their number
by a fixed \(q_K\). Write \(q_i\le q_K\) for this number.

Figure~\ref{fig:horizontal-cuts-links} shows the terminal-wall candidates
and the first-meeting decomposition of a regular piece, including its
final subpath after an earlier visit to the terminal wall.

\begin{figure}[t]
\centering
\begin{tikzpicture}[x=1cm,y=1cm,>=Latex,
  every node/.style={font=\scriptsize},
  path trace/.style={draw=black,line width=1pt},
  fine polygon/.style={draw=black!65,dashed,line width=.6pt},
  wall/.style={draw=black!60,line width=.55pt},
  wall strip/.style={fill=black!7,draw=black!22,line width=.3pt},
  support/.style={draw=black!65,densely dotted,line width=.55pt},
  cut/.style={circle,fill=black,inner sep=1.3pt}]

\begin{scope}[shift={(.42,.83)}]
  \node[font=\small,align=center] at (2.7,6.12)
    {(a) Terminal $H_m$-ball\\in $(A,B)$ coordinates};
  \draw[black!65,line width=.5pt] (0,0) rectangle (5.4,5.4);
  \foreach \v in {.5,1.3,2.1} {
    \path[wall strip] (\v-.06,-.08) rectangle (\v+.06,5.48);
    \draw[wall] (\v,-.08)--(\v,5.48);
  }
  \node[fill=white,inner sep=1pt] at (.5,5.18) {$i_1$};
  \node[fill=white,inner sep=1pt] at (1.3,5.18) {$i_2$};
  \node[fill=white,inner sep=1pt] at (2.1,5.18) {$i_m$};
  \draw[<->,black!60,line width=.35pt] (.5,4.64)--(1.3,4.64);
  \node[fill=white,inner sep=1pt] at (.9,4.64) {$J$};
  \node[fill=white,inner sep=1pt] at (1.7,4.64) {$\cdots$};
  \node[anchor=west,align=left,black!70] at (2.7,4.65)
    {$B$-walls lie in\\ $|\xi_1-i_j|\le\rho$};
  \draw[fine polygon]
    (-.15,2.12)--(.2,2.2)--(.55,2.32)--(.9,2.4)--(1.25,2.48)
    --(1.6,2.53)--(1.95,2.59)--(2.3,2.64)--(2.65,2.68)--(3,2.74);
  \draw[path trace,-{Latex[length=1.6mm]}]
    (0,2.2)--(.1,2.24)--(.2,2.2)--(.34,2.28)--(.5,2.38)--(.55,2.32)
    --(.72,2.28)--(.9,2.4)--(1.05,2.45)--(1.25,2.48)--(1.3,2.62)
    --(1.45,2.57)--(1.6,2.53)--(1.73,2.62)--(1.95,2.59)--(2.1,2.72)
    --(2.3,2.64)--(2.47,2.78)--(2.65,2.68)--(2.7,2.7);
  \foreach \px/\py in {-.15/2.12,.2/2.2,.55/2.32,.9/2.4,1.25/2.48,
                         1.6/2.53,1.95/2.59,2.3/2.64,2.65/2.68,3/2.74}
    \draw[fill=white,black!65,line width=.4pt]
      (\px-.035,\py-.035) rectangle (\px+.035,\py+.035);
  \foreach \px/\py in {.5/2.38,1.3/2.62,2.1/2.72} {
    \draw[support] (\px-.18,\py-.18) rectangle (\px+.18,\py+.18);
    \node[cut] at (\px,\py) {};
  }
  \fill (0,2.2) circle (1.35pt);
  \fill (2.7,2.7) circle (1.35pt);
  \node[anchor=east] at (-.06,2.27) {$t_-$};
  \node[anchor=east,align=right] at (-.06,1.82) {west-face\\ entry};
  \node[anchor=north] at (.5,2.08) {$c_1$};
  \node[anchor=north] at (1.3,2.32) {$c_2$};
  \node[anchor=north] at (2.1,2.42) {$c_m$};
  \node[anchor=west] at (3.11,2.74) {$P$ (fine polygon)};
  \node[anchor=west] at (2.88,2.4) {$q=\gamma(t_+)$};
  \draw[<->,black!65,line width=.35pt] (0,.4)--(2.7,.4);
  \node[fill=white,inner sep=1pt] at (1.35,.4) {$H_m$};
  \node[align=center] at (2.7,-.38)
    {Dotted boxes: candidate input supports;\\
     $J-1-2\rho>2R_0$ makes them disjoint.};
\end{scope}

\begin{scope}[shift={(6.7,1.67)}]
  \node[font=\small] at (3.25,5.28)
    {(b) A regular piece in one row band};
  \fill[black!3] (.2,1) rectangle (5.7,3);
  \draw[black!55,line width=.45pt] (0,1)--(6.05,1);
  \draw[black!55,line width=.45pt] (0,3)--(6.05,3);
  \draw[black!35,dotted,line width=.4pt] (0,2)--(6.05,2);
  \node[anchor=east] at (-.05,2) {$y_k$};
  \draw[<->,black!65,line width=.35pt] (-.17,2)--(-.17,3);
  \node[anchor=east] at (-.2,2.5) {$b$};
  \foreach \v in {.7,2.2,3.7,5.2}
    \path[wall strip] (\v-.08,.66) rectangle (\v+.08,3.4);
  \draw[wall] (.7,.66)--(.68,1.25)--(.72,1.7)--(.7,2.1)
    --(.73,2.75)--(.7,3.4);
  \draw[wall] (2.2,.66)--(2.17,1.35)--(2.22,1.95)--(2.19,2.55)
    --(2.2,3.4);
  \draw[wall] (3.7,.66)--(3.68,1.5)--(3.72,2.21)--(3.69,2.8)
    --(3.7,3.4);
  \draw[wall] (5.18,.66)--(5.24,1.25)--(5.28,1.62)--(5.21,1.9)
    --(5.16,2.28)--(5.2,2.8)--(5.18,3.4);
  \node[anchor=south] at (.7,3.42) {$W_j$};
  \node[anchor=south] at (2.2,3.42) {$W_{j+1}$};
  \node[anchor=south] at (3.7,3.42) {$W_{\ell-1}$};
  \node[anchor=south] at (5.2,3.42) {$W_\ell$};
  \node at (2.95,3.58) {$\cdots$};
  \draw[black!50,dashed,line width=.4pt] (.2,.66)--(.2,3.4);
  \draw[black!50,dashed,line width=.4pt] (5.7,.66)--(5.7,3.4);
  \draw[path trace,-{Latex[length=1.6mm]}]
    (.72,1.7)--(1.1,1.82)--(1.5,1.72)--(1.86,1.91)--(2.22,1.95)
    --(2.55,2.15)--(2.95,2.04)--(3.34,2.28)--(3.72,2.21)
    --(4.12,2.31)--(4.57,2.15)--(5.16,2.28)--(5.43,2.44)
    --(5.49,2.14)--(5.42,1.7)--(5.28,1.62);
  \foreach \px/\py in {.72/1.7,2.22/1.95,3.72/2.21,5.28/1.62}
    \node[cut] at (\px,\py) {};
  \draw[fill=white,line width=.75pt] (5.16,2.28) circle (1.7pt);
  \node[anchor=north,align=center] at (.92,.86) {initial cut};
  \draw[black!55,line width=.35pt] (.96,1.02)--(.76,1.59);
  \node[anchor=north,align=center] at (3.03,.86)
    {successive first\\ wall meetings};
  \draw[black!55,line width=.35pt] (2.66,1.04)--(2.24,1.84);
  \draw[black!55,line width=.35pt] (3.39,1.04)--(3.7,2.1);
  \node[anchor=west,align=left] at (5.95,2.75) {earlier\\ visit};
  \draw[black!55,line width=.35pt] (5.91,2.66)--(5.23,2.31);
  \node[anchor=west,align=left] at (5.95,1.36) {terminal\\ cut};
  \draw[black!55,line width=.35pt] (5.91,1.43)--(5.35,1.59);
  \draw[<->,black!65,line width=.35pt] (.2,.42)--(.7,.42);
  \draw[<->,black!65,line width=.35pt] (5.2,.42)--(5.7,.42);
  \node[anchor=north] at (.45,.37) {$b'$};
  \node[anchor=north] at (5.45,.37) {$b'$};
  \node[align=center] at (3.4,-.2)
    {terminal overshoot $\le 2r+4d\varrho<b'$};
\end{scope}
\end{tikzpicture}
\caption{Horizontal selected cuts and their use in extracting links
(schematic). In (a), the solid terminal arc starts at its last boundary hit
$\gamma(t_-)$ and ends at the mark $q=\gamma(t_+)$. The dashed fine polygon
$P$, whose fine cuts are open squares, may use partial fine intervals at
either end. The good-increment and
radius hypotheses force entry through the west face. The filled dots are
ordered first vertex meetings with the $B$-walls in the separated strips;
intermediate candidates are omitted.
Their cut cells $c_j$ have the disjoint candidate input supports represented
by the dotted boxes. In (b), a regular
piece ending on $W_\ell$ stays in $|y-y_k|\le b$ and is split at successive first meetings with
intermediate walls, most of which are omitted from the drawing. The hollow
point is an earlier visit to the terminal
wall, not the retained terminal cut. The all-pairs drawdown bound and wall
tracking give terminal overshoot at most $2r+4d\varrho<b'$, so the whole final
subpath satisfies the link window. The dashed vertical boundaries show the
outermost $b'$ margins; each unit link has the same margins.}
\label{fig:horizontal-cuts-links}
\end{figure}

\emph{Distinct columns.}
At every original cut put
\[
 v_h=\left\lfloor \xi_1(x_h)+\tfrac12\right\rfloor,
 \qquad 0\le h\le m_N.
\]
At a cut on \(W_j\), tracking gives
\(\lvert \xi_1(x_h)-j\rvert\le\rho+O(1/n)<1/2\) for large \(n\),
so \(v_h=j\). Form one integer walk by expanding each piece outside
\(\mathcal I_0\cup R\) into its \(q_i\) positive unit wall steps from
\(v_i\) to \(v_{i+1}\), and replacing each piece in
\(\mathcal I_0\cup R\) by the jump between these endpoints. Include
the two terminal jumps \(0\to v_0\) and \(v_{m_N}\to j_n\).
The walk begins at \(0\), ends at \(j_n\), and is fixed before any
whole piece is rejected by the following scan. Each jump has magnitude
\(O(K)\), including the terminal jumps. If \(q_{\rm all}\) is the
number of unit steps and \(P,D\) are the sums of the positive and
negative jump magnitudes, then
\[
 q_{\rm all}+P-D=j_n,\qquad
 P+D\le C_K\bigl(|\mathcal I_0|+|R|+1\bigr).
\]
Scan this fixed walk once, updating its running maximum at every
visited vertex. Call a unit step \(j\to j+1\) a candidate if its
endpoint is a new strict maximum and is at most \(j_n\), and let
\(q_{\rm rec}\) be the number of candidates. For each integer level
\(\ell\in\{1,\ldots,j_n\}\), the first step whose endpoint is at least
\(\ell\) is either a candidate unit step ending at \(\ell\) or a
positive jump. A positive jump is charged at most its magnitude. Hence
\[
 j_n-q_{\rm rec}\le P,\qquad
 q_{\rm all}-q_{\rm rec}
   =D-P+(j_n-q_{\rm rec})\le D.
\]
This counts every rejected unit step, including record steps beyond
\(j_n\). Let \(\mathcal C_{\rm rec}\) be the pieces outside
\(\mathcal I_0\cup R\) containing at least one rejected step, and set
\[
 \mathcal I=(\mathcal I_0\cup\mathcal C_{\rm rec})\setminus R,
 \qquad
 \mathcal G=\{0,\ldots,m_N-1\}\setminus(R\cup\mathcal I).
\]
Distinct pieces in \(\mathcal C_{\rm rec}\) contain distinct rejected
steps, so \(\lvert\mathcal C_{\rm rec}\rvert\le D\). Keep all links of
the pieces in \(\mathcal G\), without rescanning the walk. They are
candidate steps, so their columns are distinct, increasing, and in
\(\mathcal J_n\). Discarding a whole piece loses at most \(q_K\)
candidates, and therefore
\[
 |\mathcal I|\le|\mathcal I_0|+D,\qquad
 |\mathcal J_n\setminus\{j:(j,k)\in\mathcal L\}|
 \le P+q_KD.
\]
The bounds on \(P,D\) prove the stated irregular and omission bounds.
Since \(u(a)/d\) is bounded above and below, \(j_n\) is comparable
to \(N\); for sufficiently small \(\eta\), the omission bound also
gives \(\lvert\mathcal L\rvert\ge cN\) for large \(n\).

The selected rows form a chronological subsequence of the hysteresis
rows, so deleting terms cannot increase variation. Hold the old row
across each omitted gap, as in Lemma~\ref{probability:rows}. The box
in \eqref{links:long-geodesic} bounds the first selected row by
\(CN/\sigma\le j_n^2\) for large \(n\). This proves
\eqref{links:selected-rows}.

\emph{Shared data and the tunnel charge.}
After these masks are fixed, assign one datum to every original cut
\(x_h\), \(0\le h\le m_N\), retaining its selected road label and
occurrence. If at least one incident original piece belongs to
\(\mathcal G\), use that piece's row center on the cut's \(B\)-wall.
When both incident pieces belong to \(\mathcal G\), their rows agree
by the argument for closed time intervals above, and they use the same
selected wall occurrence, so they request the same center. If neither
incident piece belongs to \(\mathcal G\), use the prescribed reference
center for the cut cell and road label.
At \(h=0,m_N\), only the one existing incident piece is considered.
Use this datum for both incident pieces even if one belongs to \(R\)
or \(\mathcal I\); all incidence here is incidence in the original
list.

The masks \(R,\mathcal I_0,\mathcal C_{\rm rec}\) did not use the
eventual data, and the simultaneous \(V_{\rm cut}\) bound covers all
the prescribed choices. Thus these are the shared data used in the
telescope over the full list and the halo estimate above, which holds
for every allowed datum choice. Every piece in \(\mathcal G\) receives
its own row center at both endpoints. Any retained nonregular piece
accepting a neighboring center is already in
\(\mathcal I\), while any removed incident piece is in \(R\);
the datum choice requires no further irregular or omission charge.
For \(i\in\mathcal G\), let \(\gamma_{i,\nu}\) be its consecutive physical
subpaths from the decomposition at first meetings. Intermediate center arms
cancel, including across the retained terminal subpath, and give
\[
 \beta c_i^{\rm cost}
 =\sum_{\nu=0}^{q_i-1}
     L_{v_i+\nu,k_i}(\gamma_{i,\nu})
 \ge\sum_{\nu=0}^{q_i-1}X_{v_i+\nu,k_i}.
\]
This proves \eqref{links:unclipped-comparison}.

For each retained piece, choose the access type with its cut cells,
road labels, the two prescribed center conventions just used, and its
rounded padded domain. This type belongs to the fixed finite family,
and \(H(\gamma_i)=\beta c_i^{\rm cost}\) in
\eqref{access:center-cost}. For the tunnel comparison, choose an exterior
wall occurrence \(a_i\)
before both \(\bar x_i,x_i\), and \(b_i\) after both
\(x_{i+1},\bar x_{i+1}\), as in Lemma~\ref{access:physical}.
Additivity gives
\[
 C(a_i,\bar x_i)+\beta c_i^{\rm cost}+C(\bar x_{i+1},b_i)
 =C(a_i,x_i)+Q_n(\gamma_i)-g(x_{i+1})+g(x_i)
                                      +C(x_{i+1},b_i).
\]
The right side is the corrected cost of the physical forward
concatenation in Lemma~\ref{access:physical}. Those attachments enter
the access charge; a negative signed arm is never treated as a negative
physical traversal. Every retained piece lies in its padded \(h\)-slab
because it has no unit backward
\(h\)-excursion. Lemma~\ref{access:accounting} therefore gives
\eqref{links:tunnel-bound}. At each coarse center, sum all cut error
fields in its fixed \(O(K)\) neighborhood and over the finite cut-type
list. This fixed local field dominates the two selected endpoint errors
for every allowed piece, keeps a fixed range, and preserves the tail
exponent. The finite type and phase-template lists are those fixed in
Section~\ref{access:phases}.

Finally, one common deterministic error sequence is
\[
 \epsilon_n=C_{\rm fixed}\left(q_n+\frac{e_n}{n}+\frac nN
                         +\epsilon_n^{\rm fine}+\frac1N\right)\to0.
\]
The fixed constant absorbs \(1/\varrho,1/\sigma,1/\theta,1/c_0\)
and the bounded endpoint adjustments. The preceding cost, halo,
irregular, omission, and row estimates then have exactly the stated
bounds. Together with the uniform \(\pi_n\) above, this proves the
proposition.
\end{proof}

\paragraph{Tail loss and capping.}
Let \(\mathcal S_N\) be all finite lists on those center lattices, with
types from the fixed family, satisfying
\eqref{links:retained-geometry} and its path and multiplicity bounds.
Define \(H_N\) on the whole probability space by
\[
 H_N=\sup_{S\in\mathcal S_N}\sum_{(z,t)\in S}
       \left[Y_{z,t}\mathbf1_{\{Y_{z,t}>T\}}
          +E_z\mathbf1_{\{E_z>T\}}\right].
\]
The supremum includes the fixed type list. It is not set to zero
outside the extraction event. For a finite template domination
\(V\le\sum_{\nu=1}^{r_0}V^\nu\), use the tail comparison
\[
 V\mathbf1_{\{V>T\}}
 \le r_0\sum_{\nu=1}^{r_0}
          V^\nu\mathbf1_{\{V^\nu>T/r_0\}}.
\]
On \(\{V>T\}\), the largest summand exceeds \(T/r_0\) and the sum is
at most \(r_0\) times that largest summand. Apply this to the \(Y\)
and error fields, then sum over the fixed templates and sparse
residues. If \(1\le T<r_0\), split each summand at \(1\): values at
most \(1\) contribute at most \(C(N+1)\) by the path and multiplicity
bounds, and values above \(1\) use the tail bounds at threshold \(1\).
Since \(r_0\) is fixed and \(T<r_0\), these terms are absorbed by
\(C_F(N+1)/T+C_F\) after increasing \(C_F\). Lemma~\ref{probability:thin-tails} and
Corollary~\ref{probability:integrated-tails} give, for \(F,T\ge1\)
and uniformly in \(\mathfrak r\),
\[
 \mathbb E_{\mathfrak r}H_N
 \le \frac{CN}{F}+\frac{C_F(N+1)}{T}+C_F,
 \qquad \delta\le F^{-2}.
\]
The error fields have exponent \(m_*>4\), so their integrated bound
is included in the \(C_F/T\) term. Thus for every \(\varepsilon>0\)
one can first choose \(F\), then \(T\), and then restrict the final
cone so that
\begin{equation}
 \limsup_{n\to\infty}\sup_{\mathfrak r\in\mathfrak R_n}
       \mathbb E_{\mathfrak r}H_N/N\le\varepsilon.
 \label{links:tail-error}
\end{equation}
The test family and its collection constants remain those fixed in
the larger cone. This is an expectation bound, separate from
\(\Omega_{n,\mathfrak r}\).

For \(J,M>0\), write \(f_{J,M}(x)=\max(-J,\min(x,M))\). On the
extraction event let \(\mathcal E_J\) contain every irregular retained
piece and every regular piece containing a link below \(-J\beta\);
set \(\mathcal E_J=\varnothing\) off that event. On a nonexceptional
regular piece, replacing links by \(f_{J,M}\) only decreases the right
side of \eqref{links:unclipped-comparison}. On an exceptional regular
piece the capped sum is at most \(q_KM\), while
\eqref{links:tunnel-bound} controls the negative cost of every
exceptional retained piece. Splitting its two charges at \(T\) gives,
on \(\Omega_{n,\mathfrak r}\),
\begin{equation}
 \sum_{i\notin R}c_i^{\rm cost}
 \ge\sum_{(j,k)\in\mathcal L}f_{J,M}(X_{j,k}/\beta)
                    -(q_KM+2T)|\mathcal E_J|-H_N.
 \label{links:master}
\end{equation}
The list \(\mathcal L\) includes links in exceptional regular pieces.
Only the upper cap enters the loss coefficient.

\paragraph{Sparse negative-link flags.}
The asymmetric estimate needs a two-dimensional sparse-path bound for
\(\mathcal E_J\). A phase is shared along a wall, so column independence
alone does not give that bound. For each coarse center \(z\), let
\(\mathcal A_z\) be the fixed finite list of potential links in its
\(O(K)\) box, including its row and offset classes, and let
\[
 \Phi_z^{J}=\mathbf1_{\{\exists \ell\in\mathcal A_z:
                                      X_\ell<-J\beta\}}.
\]
Enumerate the local phase templates for the union of these links,
including both anchors and anchor means in every gauge projection.
For a fixed template \(t\), let \(\Phi_z^{J,t}\) be the same indicator
with every link evaluated under that template. At each site the
matching template reproduces every original link in \(\mathcal A_z\)
exactly, so
\[
 \Phi_z^J\le\sum_t\Phi_z^{J,t}.
\]
Because one template agrees exactly, every indicator uses the same
threshold \(J\).
The original first-moment bound gives
\(\mathbb P(X<-J\beta)\le C_s/J\). The fixed-phase first-moment bound
\eqref{links:fixed-phase-negative-moment}, followed by Markov and a
union bound over \(\mathcal A_z\), gives
\[
 \mathbb P_{\mathfrak r}(\Phi_z^{J,t}=1)\le C_K/J,\qquad J\ge1,
\]
uniformly in the law and template.

On the extraction event, every exceptional regular piece is counted
by \(\Phi_z^J\) at its center. With selected counts empty off that event,
\[
 \mathbf1_{\Omega_{n,\mathfrak r}}|\mathcal E_J|
 \le (C_K\eta+\epsilon_n)N+
   \sup_{S\in\mathcal S_N}\sum_{(z,\tau)\in S}\sum_t\Phi_z^{J,t}.
\]
When fixing the core geometry, take the sparse spacing in
Section~\ref{access:phases} large enough for the supports of
\(\mathcal A_z\), and split the center lattice into its finitely many
residues. For each residue and fixed template,
the copied indicators have disjoint product input supports; the
different template fields need not be independent. Every selected
subsequence retains the path budget and bounded multiplicity.
Lemma~\ref{probability:sparse}, summed over these fixed residues and
templates, therefore gives
\begin{equation}
 \frac{\mathbb E_{\mathfrak r}
       [\mathbf1_{\Omega_{n,\mathfrak r}}|\mathcal E_J|]}{N}
 \le C_K\eta+\epsilon_n+C_K(1+N^{-1})J^{-1/2}.
 \label{links:sparse-flags}
\end{equation}
The irregular term uses the good-event bound in the extraction
proposition; the other term bounds the supremum over all allowed
center lists before one is selected. The definition off the good
event makes all selected counts zero there.

\subsection{The two truncation estimates}

Capping makes link costs bounded, so row concentration applies.
The first estimate allows the lower cap to be chosen after a
prescribed upper cap. The second chooses equal caps from a fixed
finite menu and also controls their tail probabilities; this stronger
form will be needed for road replacement.

\begin{proposition}[Asymmetric capped-link estimate]
\label{links:asymmetric}
For each fixed \(0<M<\infty\) and \(\zeta>0\), there are
\(1\le J<\infty\), fixed geometric tolerances and \(n_0\) such that,
for every admissible \(n\ge n_0\),
\begin{equation}
 \mathbb E f_{J,M}(X/\beta)\le\zeta.
 \label{links:eq-one}
\end{equation}
The tolerances can be imposed by taking the rational \(v,w\)
sufficiently close to the endpoints of the flat. The choices are
uniform over input road laws satisfying the stated local estimates
and the sublinear upper bound for \(\beta\).
\end{proposition}

\begin{proof}
Fix an error \(\varepsilon>0\), to be chosen relative to \(\zeta\).
Choose \(F,T\) in \eqref{links:tail-error}, reserving the corresponding
upper bound for the final \(\delta\). By \eqref{links:sparse-flags},
choose \(J\) after \(T,M\) so that
\((q_KM+2T)C_KJ^{-1/2}<\varepsilon\), and then choose \(\eta\) so
that its contribution has the same bound.

Apply Lemma~\ref{probability:rows} to
\(f_{J,M}(X_{j,k}/\beta)\), using the original column independence
from Lemma~\ref{phases:locality}. Its fixed bound is \(\max(J,M)\).
The interval length \(j_n\) is comparable to \(N\); hence
\eqref{links:selected-rows} meets the lemma's row and omission
tolerances after decreasing \(\eta\) and taking \(n\) large.
Uniformly over the road law, with probability tending to one,
\[
 \left|\sum_{(j,k)\in\mathcal L}f_{J,M}(X_{j,k}/\beta)
       -|\mathcal L|\mathbb E f_{J,M}(X/\beta)\right|
 \le\varepsilon N
\]
on the extraction event. The supremum in the row lemma is taken before
the rows and omissions are selected.

If \eqref{links:eq-one} failed uniformly, choose an offending sequence
of scales and road laws. On the geometric and row events the capped
sum is at least \(c\zeta N-\varepsilon N\), while
\eqref{links:retained-cost} bounds the left side of
\eqref{links:master} by \(\epsilon_nN\). Since \(\mathcal E_J\) is
empty off the good event, \eqref{links:sparse-flags} and
\eqref{links:tail-error} bound the expected nonnegative losses in
\eqref{links:master} by \(C\varepsilon N\) for all large \(n\).
By Markov's inequality, they are at most \(4C\varepsilon N\) with
probability at least \(3/4\). This event intersects the geometric and row events.
Choosing \(\varepsilon<c\zeta/(10C+10)\) contradicts
\eqref{links:master}. Every threshold and geometric tolerance was
fixed before \(n\), proving the asserted uniformity.
\end{proof}

The same estimate also bounds the untruncated positive-tail
probability. Since \(f_{J,M}(X/\beta)\) is at least \(M\) on
\(\{X>M\beta\}\) and at least \(-(X/\beta)^-\) elsewhere,
\begin{equation}
 \mathbb P(X>M\beta)\le(C_s+\zeta)/M.
 \label{links:positive-probability}
\end{equation}
This uses no assumption that the positive tails are uniformly integrable.

\begin{proposition}[Matched caps in a finite dyadic range]
\label{links:matched}
Let \(X^{(1)},\ldots,X^{(r)}\) be finitely many required row-link
templates, and let \(U_1,\ldots,U_p\) be additional nonnegative local
variables with a uniform bound on the sum of their first moments.
For every \(\zeta>0\) and every prescribed lower bound
\(M_{\rm prescribed}\), there is a fixed finite dyadic set
\(\mathcal M=\{M_0,2M_0,\ldots,2^{H-1}M_0\}\), with
\(M_0\ge M_{\rm prescribed}\), and fixed geometric tolerances, such
that for every sufficiently large \(n\) some deterministic law-dependent
\(M\in\mathcal M\) satisfies, simultaneously,
\begin{align}
 \mathbb E f_{M,M}(X^{(a)}/\beta)&\le\zeta
       &&(1\le a\le r),\label{links:matched-mean}\\
 M\left[\sum_{a=1}^r\mathbb P((X^{(a)}/\beta)^->M)
       +\sum_{b=1}^p\mathbb P(U_b>M)\right]&\le\zeta.
 \label{links:matched-tail}
\end{align}
In particular \(U_1\) may be the positive part of the normalized
primary increment used in the road improvement.
\end{proposition}

\begin{proof}
Fix an error \(\varepsilon>0\). Choose \(F,T\) as in
\eqref{links:tail-error}, uniformly over the finite type family,
and reserve the corresponding upper bound for the final \(\delta\).
Choose \(M_0\ge M_{\rm prescribed}\) so large that
\(2TC_s/M_0<\varepsilon\). Let \(A\) bound the sum of the first
moments of \((X^{(a)}/\beta)^-\) and the \(U_b\). Choose \(H\) so
large that \(2A/H<\varepsilon\). Lemma~\ref{probability:dyadic}
ensures that for each law some level in the fixed menu satisfies
\begin{equation}
 M\sum_a\mathbb P(X^{(a)}<-M\beta)
       +M\sum_b\mathbb P(U_b>M)\le\varepsilon.
 \label{links:good-dyadic-level}
\end{equation}
We fix all concentration and geometric tolerances for the whole menu
before selecting this law-dependent level.

For equal caps, the square-root probability in
\eqref{links:sparse-flags} would be multiplied by the cap. Instead use
the original row-link fields and their column independence. For every
\(M\in\mathcal M\) and every link template, apply
Lemma~\ref{probability:rows} separately to
\(\mathbf1_{\{X_{j,k}^{(a)}<-M\beta\}}\) and to the capped costs.
Take a finite union of these conclusions, using
\(M_{\max}=2^{H-1}M_0\) as a common bound. The resulting event has
probability tending to one uniformly over the law, and controls all
candidate levels before one is selected.

On this event and the corresponding extraction events, the selected
negative-link indicators have frequency at most their marginal
probability plus a fixed error \(\nu\). Columns are distinct, so each
exceptional regular piece is counted by at least one such indicator.
For any candidate \(M\), the regular exceptions in
\eqref{links:master} with \(J=M\) therefore cost at most
\[
 C_K(M+T)N\left[\sum_a\mathbb P(X^{(a)}<-M\beta)+\nu\right].
\]
The irregular contribution is at most
\((q_KM+2T)(C_K\eta+\epsilon_n)N\).
After fixing \(M_{\max}\), choose \(\nu\), the row and omission
tolerances, and \(\eta\) so that their contributions are at most
\(\varepsilon N\). The term containing \(\epsilon_n\) then has the
same bound for large \(n\). Select a level satisfying
\eqref{links:good-dyadic-level}. Its \(M\)-weighted term is
\(O(\varepsilon N)\); its \(T\)-weighted term is small by the same
bound or by \(\mathbb P(X<-M\beta)\le C_s/M_0\).

The capped row sums at this level are compared with their means by
the same simultaneous concentration event. The extraction gives
\(|\mathcal L|\ge cN\) and retained cost at most \(\epsilon_nN\)
for each template. Markov applied to the finite sum of their \(H_N\)'s
supplies a loss event of fixed positive probability, which intersects
all the events whose probabilities tend to one. The contradiction in the asymmetric proof,
with \(\varepsilon\) sufficiently small relative to \(\zeta\),
therefore proves \eqref{links:matched-mean} for every template.
Equation~\eqref{links:good-dyadic-level} gives
\eqref{links:matched-tail}. The same cap is used for both conclusions;
no expectation bound on a discarded heavy negative tail is required.
\end{proof}
